%% kapa: exam=1 round=22 year=2011 period=1 form=A subject=economics subject_ko=경제원론 seq=7 number=47
%% source: Q-Net 감정평가사 기출문제 (cst00309) artlSeq=5036790 file=extracted/1008/5036790/2011년 제22회 감정평가사 1차시험 기출문제/1교시.hwp
%% conversion: kordoc 4.24.51 -> tools/gen1.py (auto); review=proofread-sonnet; flags=none
\begin{question}{47}
다음 ( )안에 들어갈 내용이 순서대로 올바른 것은?
\begin{framedtext}
J-curve 효과는 ‘환율이 ( ㄱ )하면 한국의 경상수지가 초기에는 ( ㄴ )되고 시간이 경과된 후에는(도) ( ㄷ )되는 효과가 나타나는 것’을 의미한다. (단, 환율은 미국 달러에 대한 원화의 환율: ￦/＄, 양국의 물가수준은 불변)
\end{framedtext}
\begin{choices}
  \item ㄱ. 상승, ㄴ. 악화, ㄷ. 개선
  \item ㄱ. 상승, ㄴ. 개선, ㄷ. 개선
  \item ㄱ. 상승, ㄴ. 악화, ㄷ. 악화
  \item ㄱ. 하락, ㄴ. 악화, ㄷ. 개선
  \item ㄱ. 하락, ㄴ. 악화, ㄷ. 불변
\end{choices}
\end{question}
\answer{1}
